\documentclass[12pt, letterpaper]{article} 

\usepackage{amsmath, amssymb, amsthm, enumerate, graphicx}

\title{MATH 189 HW 2}
\author{Dani Nguyen}
\date{Winter 2025}

\begin{document}
\maketitle
\graphicspath{ {./} }

\noindent \textbf{Question 1}
\begin{enumerate}[a.]
    \item Since IQ and GPA is fixed, we will denote those values as a constant $C$
    \begin{align*}
        Y &= C + 35X_3 - 10X_5 \\
        &= C + \mathbb{I}\{\text{has college}\} (35-10X_5)
    \end{align*}
    When fixing IQ and GPA, the model predicts that college graduates will earn less 
    if their GPA is higher than 3.5. Therefore, option iii. is correct.
    \item For a college graduate with IQ 110 and GPA 4.0,
    \begin{align*}
        Y &= 50 + 20(4.0) + 0.07(110) + 35(1) + 0.01(4.0 \cdot 110) - 10(4.0\cdot1) \\
        &= 137.1 \text{ thousand dollars}
    \end{align*}
    \item False, because GPA and IQ are on average very large numbers, especially 
        when multiplied, that even a low coefficient can lead to large effects in
        salary. This does not mean that there is little evidence of an interaction 
        effect, one would need to look at the $p$-values of the interaction term to 
        know whether there exists an interaction effect. 
\end{enumerate}
\newpage
\noindent \textbf{Question 2} \\
\begin{align*}
    \hat y_i &= x_i \hat \beta \\
    &= x_i \frac{\sum_{i=1}^{n}x_iy_i}{\sum_{i'=1}^{n}x_{i'}^2} \\
    &= \frac{\sum_{i=1}^{n}x_i^2 y_i}{\sum_{i'=1}^{n}x_{i'}^2} \\
    &= \sum_{i=1}^{n} \frac{x_i^2}{\sum_{i'=1}^{n}x_{i'}^2} y_i
\end{align*}
Let $a_i = \frac{x_i^2}{\sum_{i'=1}^{n}x_{i'}^2}$
\begin{align*}
    \hat y_i = \sum_{i=1}^{n}a_iy_i
\end{align*}

\noindent \textbf{Question 3} \\
Recall that $\beta_0 = \bar y + \beta_1 \bar x$. This can be rearranged to 
$\bar y = \beta_0 - \beta_1\bar x$. This means that the equation for simple linear 
regression will always pass through the point $(\bar x, \bar y)$.

\noindent \textbf{Question 4}
\begin{align*}
    R^2 &= \frac{TSS-RSS}{TSS} \\
    &= \sum_{i=1}^{n}(y_i)^2
    Cor(x,y) &= \frac{\sum_i (x_i-\bar{x})(y_i - \bar{y})}{\sqrt{\sum_i(x_i - \bar{x})^2}\sqrt{\sum_i(y_i - \bar{y})^2}}
\end{align*}

\end{document}